\section{Framework}
\label{sec:framework}

\subsection{Setup}

Fix a task distribution $\mathcal{T}$ and a generator $\pi$ (a language model
with a decoding temperature). For a task $t$, drawing a sample $y \sim \pi(t)$
yields a correct answer with probability $\pone(t)$, and we write
$\pone = \mathbb{E}_t[\pone(t)]$ for the single-sample accuracy. A
\emph{$k$-agent ensemble} draws $y_1,\dots,y_k \sim \pi(t)$ independently and
applies a \emph{selector} $\sigma$ that returns one index. We write
$A_\sigma(k)$ for the resulting accuracy.

Two reference selectors bound the rest. The \emph{oracle} selector returns a
correct candidate whenever one exists, so its accuracy is the coverage
\begin{equation}
\cov(k) \;=\; \mathbb{E}_t\bigl[\Pr(\exists\, i \le k : y_i \text{ correct})\bigr],
\end{equation}
which we estimate with the unbiased $\mathrm{pass}@k$ estimator of
\citet{chen2021codex} from a pool of $n \ge k$ samples. The
\emph{uninformative} selector chooses an index independently of correctness;
its accuracy is $\pone$ for every $k$, since each candidate is marginally
correct with probability $\pone$.

\subsection{The coverage--selection decomposition}

\begin{definition}[Selection efficiency]
For a selector $\sigma$ with $\cov(k) \neq \pone$,
\begin{equation}
\eff_\sigma(k) \;=\; \frac{A_\sigma(k) - \pone}{\cov(k) - \pone}.
\label{eq:eta}
\end{equation}
\end{definition}

\begin{proposition}[Decomposition]
\label{prop:decomp}
For any selector, $A_\sigma(k) = \pone + \eff_\sigma(k)(\cov(k)-\pone)$.
An uninformative selector has $\eff_\sigma(k)=0$ and the oracle has
$\eff_\sigma(k)=1$. Every selector satisfies $\eff_\sigma(k) \le 1$;
$\eff_\sigma(k) < 0$ is attainable and indicates a selector anti-correlated
with correctness.
\end{proposition}

\begin{proof}
The identity is Equation~\eqref{eq:eta} rearranged. For the uninformative
selector, selection is independent of correctness, so the chosen candidate is
correct with its marginal probability $\pone$, giving $\eff=0$. For the oracle
$A_\sigma(k)=\cov(k)$ gives $\eff=1$. No selector can exceed $\cov(k)$, since it
returns one of the generated candidates and is correct only if that candidate
is; hence $\eff_\sigma(k)\le 1$. Nothing prevents $A_\sigma(k) < \pone$---a
selector that systematically prefers confident errors---so $\eff_\sigma(k)$ is
not bounded below by zero.
\end{proof}

Proposition~\ref{prop:decomp} is elementary, and deliberately so: its value is
that it names the two factors that any multi-agent gain must pass through, and
both are separately measurable. It also makes the failure modes explicit.
A selector cannot rescue a pool with no headroom ($\cov(k) \approx \pone$), and
headroom is worthless without a selector that can find it ($\eff \approx 0$).
The possibility $\eff < 0$ is not hypothetical; we observe it
(Section~\ref{sec:results}).

\subsection{From pairwise discrimination to listwise selection}

The decomposition is useless prospectively unless $\eff$ can be estimated
without building the system. We estimate it from a cheaper probe.

\begin{definition}[Generator--verifier gap]
Present the model with a task and two candidate answers, one correct and one
incorrect, in randomised order, and ask which is correct. Let $q$ be the
probability it answers correctly. The generator--verifier gap is
$\gap = 2q - 1 \in [-1, 1]$.
\end{definition}

$\gap=0$ means the model cannot distinguish its own correct answers from its
incorrect ones; $\gap=1$ means it always can; $\gap<0$ means it is
systematically misled. Crucially $\gap$ is a \emph{pairwise} quantity, while a
$k$-agent ensemble makes a \emph{listwise} choice. We bridge the two with a
minimal latent-score model.

\begin{definition}[Latent-score selector]
\label{def:latent}
The verifier assigns candidate $i$ a score $s_i$, drawn independently from
$\mathcal{N}(\delta,1)$ if $y_i$ is correct and $\mathcal{N}(0,1)$ otherwise,
and selects $\arg\max_i s_i$.
\end{definition}

\begin{proposition}[Identification and prediction]
\label{prop:latent}
Under Definition~\ref{def:latent}, the pairwise preference probability is
$q = \Phi(\delta/\sqrt{2})$, so an observed gap identifies the separation
\begin{equation}
\delta = \sqrt{2}\,\Phi^{-1}\!\left(\tfrac{1+\gap}{2}\right)
\label{eq:delta}
\end{equation}
with no free parameters. If $c$ of the $k$ candidates are correct, the
probability of selecting a correct one is
\begin{equation}
P(c,k,\delta) \;=\; \Pr\Bigl(\max_{j\le c} X_j > \max_{j \le k-c} Z_j\Bigr),
\quad X_j \sim \mathcal{N}(\delta,1),\; Z_j \sim \mathcal{N}(0,1),
\label{eq:listwise}
\end{equation}
and the predicted selector accuracy is
$\hat{A}_\sigma(k) = \mathbb{E}_{c}\bigl[P(c,k,\delta)\bigr]$, with the
expectation taken over the observed distribution of $c$. The predicted
selection efficiency follows from Equation~\eqref{eq:eta}.
\end{proposition}

\begin{proof}
$s_{\text{correct}} - s_{\text{incorrect}} \sim \mathcal{N}(\delta, 2)$, so
$q = \Pr(s_{\text{correct}} > s_{\text{incorrect}}) = \Phi(\delta/\sqrt{2})$,
which inverts to Equation~\eqref{eq:delta} since $\Phi$ is strictly increasing.
Equation~\eqref{eq:listwise} is the definition of the argmax event conditional
on $c$; $P(0,k,\delta)=0$ and $P(k,k,\delta)=1$. Averaging over $c$ gives
$\hat{A}_\sigma(k)$.
\end{proof}

Two things make Proposition~\ref{prop:latent} usable. First, $\gap$ and the
distribution of $c$ are both estimable from a small labelled calibration set by
sampling $k$ answers per task---no multi-agent system need be built. Second,
nothing is fitted: Equation~\eqref{eq:delta} pins $\delta$, and
Equation~\eqref{eq:listwise} is evaluated by simulation. The model is
intentionally the simplest that respects the pairwise/listwise distinction; it
assumes scores are conditionally independent and equally dispersed, and
Section~\ref{sec:results} reports where those assumptions hold and where they
fail.

\subsection{Aggregation without a verifier}

Self-consistency uses no verifier, so $\gap$ cannot explain it. It relies
instead on the sampling distribution concentrating on the truth.

\begin{definition}[Plurality alignment]
Let $\hat{y}$ be the modal answer among $k$ samples. Plurality alignment is
$\plur = \Pr(\hat{y} \text{ correct}) - \pone$, the accuracy advantage of the
mode over a single sample.
\end{definition}

$\plur$ is also estimable from the same calibration pool at no extra cost, and
it is the natural predictor of self-consistency's lift. Separating $\gap$ from
$\plur$ matters empirically: we find families where $\plur > 0$ while $\gap
\le 0$, so a system can benefit from voting while the very same model is
useless as a judge.

\subsection{Token-matched lift and the decision rule}
\label{sec:matched}

Let $\mathcal{C}(\cdot)$ denote total tokens consumed per task, prompt plus
completion. A $k$-agent ensemble costs roughly $k$ times a single sample. We
therefore compare it not against one sample but against the \emph{single-agent
frontier} $a(\cdot)$: the accuracy a single agent reaches when given the same
budget to spend serially, by generating and then revising its own answer.

\begin{definition}[Token-matched lift]
For an architecture $\mathcal{A}$ with cost $\mathcal{C}(\mathcal{A})$,
\begin{equation}
\lift(\mathcal{A}) \;=\; A(\mathcal{A}) \;-\; a\bigl(\mathcal{C}(\mathcal{A})\bigr).
\end{equation}
\end{definition}

$\lift > 0$ is the claim a multi-agent architecture must support: that spending
the budget across coordinating agents beats spending it inside one. Combining
with Proposition~\ref{prop:decomp} gives the rule we evaluate prospectively:

\begin{equation}
\textbf{build the ensemble} \iff
\underbrace{\hat{\eff}_\sigma(k)\bigl(\hat{\cov}(k) - \hat{\pone}\bigr)}_{\text{predicted from calibration}}
\;>\;
\underbrace{a\bigl(\mathcal{C}(\mathcal{A})\bigr) - \pone}_{\text{sequential gain}} .
\label{eq:rule}
\end{equation}

Every quantity on the left is measurable in advance from
Proposition~\ref{prop:latent} and repeated sampling; the right-hand side
requires only running one agent at a larger budget. Section~\ref{sec:results}
tests both the numerical prediction and the binary decision it implies.
